% ============ 04. Datos utilizados: suficiencia y corrección ============

\section{Datos utilizados: suficiencia y corrección}

\subsection{Inventario de fuentes de datos}

\begingroup
\small
\setlength{\tabcolsep}{3pt}
\begin{longtable}{>{\raggedright\arraybackslash}p{3.0cm}>{\raggedright\arraybackslash}p{4.9cm}>{\raggedright\arraybackslash}p{4.1cm}>{\raggedright\arraybackslash}p{2.6cm}}
\caption{Inventario de fuentes de datos del repositorio}
\label{tab:fuentes}\\
\toprule
\textbf{Fuente} & \textbf{Descripción} & \textbf{Volumen} & \textbf{Estado en el clon} \\
\midrule
\endfirsthead
\toprule
\textbf{Fuente} & \textbf{Descripción} & \textbf{Volumen} & \textbf{Estado en el clon} \\
\midrule
\endhead
Padrón de investigadores & Dataset Socrata \texttt{bqtm-4y2h}: investigadores reconocidos por convocatoria, 6 convocatorias (2013--2021), 30 variables &
77\,237 filas · 30\,086 únicos &
Consolidado XLSX/CSV versionado en \texttt{datos/tarea\_join/} \\
Producción de grupos & Dataset Socrata \texttt{33dq-ab5a}: una fila por producto-autor &
3\,166\,629 filas · 77\,401 autores &
\textbf{No presente} (gitignored, ~1.2\,GB) \\
Catálogo & \texttt{datos/catalogo.yaml}: metadatos, diccionario de datos, llaves de cruce &
308 líneas &
Versionado \\
Gran tabla (sin ID) & Tabla analítica sin identificadores personales &
— (mencionada en \texttt{informe\_final.md}) &
\textbf{No versionada} (solo referencias en notebooks legacy) \\
\bottomrule
\end{longtable}
\endgroup

\subsection{¿Los datos eran suficientes para el objetivo?}

\begin{tcolorbox}[hallazgo]
\textbf{Sí, los datos son suficientes} para el objetivo descriptivo-exploratorio
planteado: 6 convocatorias permiten análisis longitudinal de retención y
transición; 30\,086 investigadores únicos dan potencia suficiente para
desagregaciones por área, género y región; y el cruce con 3\,166\,629 productos
permite analizar productividad. La cobertura temporal (2013--2021) es adecuada
para estudiar una década de reconocimiento científico en Colombia.
\end{tcolorbox}

Limitaciones de suficiencia (detalladas en la sección de alcance y riesgos):

\begin{itemize}
  \item \textbf{Subgrupos pequeños}: raizales, palenqueros, Rrom y personas con
        discapacidad tienen conteos reducidos $\rightarrow$ estimaciones inestables,
        intervalos amplios.
  \item \textbf{Ventanas irregulares}: los intervalos entre convocatorias no son
        constantes (2013, 2014, 2015, 2017, 2019, 2021), lo que complica la
        comparación de tasas de transición entre periodos.
  \item \textbf{Ausencia de convocatorias recientes}: no se incluye 2023; el
        README lo señala como pendiente (Sprint 1).
  \item \textbf{Producción no reproducible localmente}: el dataset de producción
        no está versionado; su ingesta depende de la API Socrata (~1.2\,GB).
\end{itemize}

\subsection{¿Los datos eran los correctos para el objetivo?}

\begin{tcolorbox}[hallazgo]
\textbf{Los datos son los correctos} para el objetivo planteado: el padrón de
investigadores reconocidos es la fuente oficial del sistema de reconocimiento de
MinCiencias y el dataset de producción es el complemento natural para medir
outputs. La elección de fuentes de datos abiertos (datos.gov.co, Socrata) es
apropiada y verificable.
\end{tcolorbox}

Problemas de corrección detectados en la auditoría:

\begin{enumerate}
  \item \celdanok{Codificación (mojibake)}: el CSV canónico
        \texttt{investigadores\_consolidado.csv} presenta caracteres corruptos
        (\texttt{Física} por \emph{Física}, \texttt{Bogotá} por \emph{Bogotá}),
        un desajuste UTF-8/Latin-1 en la cadena de exportación Socrata $\rightarrow$
        CSV. El XLSX no se ve afectado.
  \item \celdanok{Atípicos de edad}: 22 registros con edad $>100$ años (máximo 956),
        tratados por eliminación documentada; inconsistencia documental 10 vs.\ 22.
  \item \celdanok{Edad fraccionaria}: la variable \texttt{EDAD\_ANOS\_PR} contiene
        valores decimales (p.~ej. 54.48), lo que sugiere edad calculada con
        precisión de días y posterior redondeo inconsistente.
  \item \celdanok{Faltantes estructurales}: \texttt{NME\_DEPARTAMENTO\_*}
        ``NO DISPONIBLE'' 5.2\%; género ``NO REPORTADO'' 2.7\%; cobertura de
        etnia/discapacidad/conflicto 0\% antes de 2021.
  \item \celdanok{Etiquetado de convocatoria}: la convocatoria 833 (2018) registra
        \texttt{ANO\_CONVO = 2019} (fecha de publicación de resultados); está
        documentado en el catálogo, pero puede inducir a error si no se controla.
  \item \celdanok{Consolidado versionado incompleto (3 de 6 convocatorias)}: tanto
        \texttt{datos/tarea\_join/investigadores\_consolidado.xlsx} como el CSV
        paralelo contienen \textbf{solo 50\,891 filas y 3 convocatorias}
        (2017: 13\,001, 2019: 16\,796, 2021: 21\,094), frente a las 77\,237 filas
        y 6 convocatorias (2013--2021) declaradas en el README y el catálogo.
        Las evidencias (p.~ej. matrices de transición 2013--2014) sí cubren las
        6 convocatorias, lo que indica que el artefacto versionado es un
        \emph{snapshot} de la fase inicial de 3 convocatorias que nunca se
        actualizó. Consecuencia: \textbf{un clon fresco que ejecute el pipeline
        no reproduce los hallazgos del README} (faltan 2013, 2014 y 2015);
        la fuente completa vive solo en \texttt{datos/raw} (gitignored).
        Hallazgo verificado con conteo real (pandas 3.0) sobre ambos archivos.
  \item \celdanok{Captura selectiva de co-filiación}: el campo \texttt{INST\_FILIA}
        solo capturó múltiples instituciones en 2019 (569 casos), por lo que el
        análisis de redes es un corte transversal, no una serie.
  \item \celdanok{Identificadores}: se asumió la unicidad de \texttt{ID\_PERSONA\_PR}
        con una heurística de coherencia de edad ($\pm$4 años) y género; no hay
        validación formal contra fuentes externas.
\end{enumerate}

\subsection{Evaluación de la calidad de datos por dimensión (Wang \& Strong)}

\begin{table}[htbp]
\centering
\caption{Estado de la calidad de datos por dimensión (Wang \& Strong, 1996)}
\label{tab:calidad_dimensiones}
\small
\begin{tabularx}{\linewidth}{p{3.6cm}p{7.2cm}p{3.6cm}}
\toprule
\textbf{Dimensión} & \textbf{Observación} & \textbf{Estado} \\
\midrule
Exactitud & Cruce con fuentes oficiales no realizado; mojibake en CSV &
\celdanok{Parcial} \\
Completitud & Mapeada por convocatoria (heatmaps); faltantes estructurales declarados &
\celdaok{Adecuada} \\
Consistencia & Encodings inconsistentes; duplicidad de artefactos; etiquetas normalizadas en transformación &
\celdanok{Parcial} \\
Unicidad & Verificación de IDs por coherencia de edad y género; sin formalizar &
\celdaok{Adecuada} \\
Validez & Atípicos de edad documentados y tratados con justificación &
\celdaok{Adecuada} \\
Oportunidad & Datos de 2021 con pull 2026; convocatoria 2023 pendiente &
\celdanok{Parcial} \\
\bottomrule
\end{tabularx}
\end{table}

\subsection{Conclusión de la sección}

Los datos son \textbf{correctos y suficientes} para el objetivo; sin embargo, la
auditoría detecta \textbf{problemas de calidad corregibles} (codificación,
faltantes, atípicos, captura selectiva) que deben documentarse explícitamente en
cada entregable y corregirse en el pipeline de transformación (normalización
UTF-8, contrato de esquema, imputaciones documentadas).
